%----------------------------------------------------------
% 제19장. 제조/산업 AI 거버넌스
%----------------------------------------------------------

\chapter{제조/산업 AI 거버넌스}
\label{ch:industrial_governance}

제조 및 산업 현장의 AI는\cite{li2023trustworthy} 다른 어떤 영역보다 \textbf{물리적 안전}과 직결된다. 금융 AI의 오류가 경제적 손실을 야기한다면, 산업 AI의 오류는 작업자의 생명과 신체에 직접적 위협을 가할 수 있다. 이러한 특수성으로 인해 산업 AI 거버넌스는 \textbf{기능안전(Functional Safety)} 개념을 핵심 축으로 삼아야 한다.

산업 AI의 또 다른 특수성은 \textbf{OT(Operational Technology)}와 \textbf{IT(Information Technology)}의 융합 환경에서 작동한다는 점이다. 전통적으로 폐쇄망에서 운영되던 제어 시스템이 AI 도입과 함께 IT 네트워크와 연결되면서, 사이버 보안 위협이 물리적 안전 위협\cite{lee2018cyber}으로 전이될 수 있는 Risk가 등장했다.

본 장에서는 산업 AI의 특수성을 분석하고, IEC 61508, ISO 13849, UL 4600 등 기능안전 표준의 AI 적용 방안을 제시하며, 예지보전\cite{wuest2016machine}, 품질검사, 자율로봇, 디지털 트윈\cite{tao2018digital}, 해양/조선 등 주요 산업 AI\cite{wang2018towards} 응용 영역의 거버넌스 사례를 심층적으로 다룬다.

%----------------------------------------------------------
\section{산업 AI의 특수성}
\label{sec:19.1}

\subsection{안전 중심(Safety-Critical) 시스템}
\label{sec:19.1.1}

산업 AI는 대부분 시스템의 실패가 인명 피해나 환경 오염으로 이어지는 \textbf{안전 중심 시스템}의 일부로 작동한다. 따라서 실시간성, 확정성(Determinism), 가용성이 극도로 중요하다.

\begin{table}[htbp]
\centering
\caption{산업 AI의 안전 중심 특성}
\label{tab:19_1}
\begin{tabularx}{\textwidth}{p{2cm}Xp{5cm}}
\toprule
\textbf{특성} & \textbf{내용} & \textbf{거버넌스 시사점} \\
\midrule
실시간성 & 밀리초 단위 의사결정 요구 & 지연 시간 보장, 실시간 OS 연계 \\
확정성 & 동일 입력에 동일 출력 보장 & 비결정론적 모델의 한계 관리 \\
가용성 & 24/7 무중단 운영\cite{klaise2021monitoring} & 이중화, 페일세이프(Fail-safe) 설계 \\
추적가능성 & 모든 결정 사후 추적 & 블랙박스 Logging, 결정 이력 관리 \\
\bottomrule
\end{tabularx}
\end{table}

\paragraph{안전성 Integrity Level (SIL)}

IEC 61508은 시스템의 안전 무결성을 \textbf{SIL 1~4}로 분류한다 \cite{iec61508}. AI 모델은 본질적으로 통계적 불투명성을 가지므로\cite{arrieta2020explainable}, 산업 AI 거버넌스는 AI를 직접 제어가 아닌 \textbf{보조/감시(Advisory)} 역할로 한정하고, 하위 계층에 검증된 하드웨어 안전 장치를 두는 \textbf{다층 방어(Defense in Depth)} 전략을 취한다.

\begin{lstlisting}[language=Python, caption={Industrial AI Safety Integrity (SIL) Assessment and ODD Constraint Logic}, label={lst:sil_eval}]
class IndustrialAISafetyAssessor:
 """AI Safety Integrity Assessment based on IEC 61508"""
 def __init__(self):
 # Max allowed SIL by AI involvement level (Limited to SIL 1 for direct control)
 self.ai_max_sil = {'direct': 1, 'advisory': 2, 'monitoring': 3}

 def assess_ai_function(self, involvement, target_sil, diagnostic_coverage):
 max_allowed = self.ai_max_sil.get(involvement, 0)

 if target_sil > max_allowed:
 return "Non-compliant: SIL > 1 not allowed for direct AI control. Hardware interlocking required."

 if diagnostic_coverage < 0.90:
 return "Supplement required: 90%+ diagnostic coverage needed for SIL 2 or higher."

 return "Compliant: Approval for operation within designed ODD (Operation Design Domain)."
\end{lstlisting}

\subsection{OT-IT 융합과 보안 거버넌스}
\label{sec:19.1.2}

산업 AI는 퍼듀 모델(Purdue Model) 기반의 계층형 아키텍처에서 작동한다. 보안 거버넌스는 IT 계층의 공격(Attack)이 OT 계층으로 전이되지 않도록 \textbf{산업 DMZ}를 통한 망 분리와 \textbf{데이터 다이오드(Data Diode)}를 활용해야 한다.

\paragraph{퍼듀 모델과 산업 AI 배치}

퍼듀 모델은 산업 제어 시스템을 5개 수준(Level 0~4)으로 계층화하는 참조 아키텍처이다. 각 수준에서 AI의 역할과 보안 요구사항이 다르다.

\begin{table}[htbp]
\centering
\caption{퍼듀 모델 수준별 AI 배치와 보안 요건}
\label{tab:19_purdue}
\begin{tabularx}{\textwidth}{p{1.5cm}p{2.5cm}Xp{3.5cm}}
\toprule
\textbf{수준} & \textbf{계층} & \textbf{AI 역할} & \textbf{보안 요건} \\
\midrule
Level 0 & 물리적 프로세스 & AI 직접 제어 불가 (센서/액추에이터) & 물리적 격리, 하드와이어드 안전 장치 \\
Level 1 & 기본 제어 (PLC) & AI 보조 제어 (SIL 1 한정) & 전용 프로토콜, 화이트리스트 \\
Level 2 & 감시 제어 (HMI/SCADA) & AI 모니터링, 이상 탐지 & 접근 제어, 감사 로그 \\
Level 3 & 운영 관리 (MES) & 예지보전, 품질 예측, 최적화 & 산업 DMZ 경유 접근 \\
Level 3.5 & 산업 DMZ & 데이터 다이오드, 단방향 게이트웨이 & 양방향 통신 금지 \\
Level 4 & 기업 IT (ERP) & 데이터 분석, 대시보드, 보고 & 일반 IT 보안 정책 \\
\bottomrule
\end{tabularx}
\end{table}

\paragraph{IEC 62443 산업 사이버 보안 표준}

IEC 62443은 산업 자동화 및 제어 시스템(IACS)의 사이버 보안을 위한 국제 표준이다. AI가 산업 시스템에 도입될 때 반드시 준수해야 하는 보안 프레임워크를 제공한다.

\begin{table}[htbp]
\centering
\caption{IEC 62443 보안 수준(Security Level)과 AI 시스템 요구사항}
\label{tab:19_iec62443}
\begin{tabularx}{\textwidth}{p{1.5cm}p{3cm}Xp{3cm}}
\toprule
\textbf{SL} & \textbf{위협 수준} & \textbf{대응 능력} & \textbf{AI 시스템 요건} \\
\midrule
SL 1 & 우연한 비인가 접근 & 기본 보안 & 기본 접근 제어, 로그 \\
SL 2 & 의도적 공격 (낮은 역량) & 방어적 보안 & 암호화, 무결성 검증 \\
SL 3 & 정교한 공격 (높은 역량) & 심층 방어 & 적대적 공격 방어, 이상 탐지 \\
SL 4 & 국가 수준 위협 & 최고 수준 방어 & AI 사용 제한, 하드웨어 격리 \\
\bottomrule
\end{tabularx}
\end{table}

\paragraph{데이터 다이오드와 단방향 게이트웨이}

데이터 다이오드(Data Diode)는 물리적으로 단방향 통신만 허용하는 보안 장치이다. OT 영역의 센서 데이터를 IT 영역의 AI 시스템으로 전송할 때, 역방향 통신(IT에서 OT로)을 물리적으로 차단하여 외부 공격이 제어 시스템에 도달하는 것을 원천 방지한다. 산업 AI에서 데이터 다이오드는 다음과 같이 활용된다:

\begin{itemize}
    \item \textbf{센서 데이터 수집}: OT 센서 데이터를 단방향으로 AI 분석 서버에 전송
    \item \textbf{예지보전 분석}: AI 분석 결과를 별도 경로(산업 DMZ 경유)로 운영자에게 전달
    \item \textbf{로그 수집}: 제어 시스템 로그를 단방향으로 보안 모니터링 시스템에 전송
\end{itemize}

%----------------------------------------------------------
\section{산업 AI 표준과 인증}
\label{sec:19.2}

\subsection{IEC 61508 SIL 상세}
\label{sec:19.2.1}

IEC 61508은 전기/전자/프로그래머블 전자 안전 관련 시스템(E/E/PE Safety-related Systems)의 기능안전 기본 표준이다. AI가 산업 안전 시스템에 적용될 때 이 표준의 요구사항을 이해하고 준수하는 것이 거버넌스의 출발점이다.

\paragraph{SIL 1-4 수준별 요구사항}

SIL(Safety Integrity Level)은 안전 기능이 요구 시 올바르게 작동할 확률을 나타내는 지표이다. SIL 1이 가장 낮고 SIL 4가 가장 높은 안전 무결성을 요구한다.

\begin{table}[htbp]
\centering
\caption{IEC 61508 SIL 수준별 요구사항 상세}
\label{tab:19_sil_detail}
\begin{tabularx}{\textwidth}{p{1cm}p{2.5cm}p{2.5cm}Xp{2.5cm}}
\toprule
\textbf{SIL} & \textbf{목표 고장 확률 (PFD$_\text{avg}$)} & \textbf{위험측 고장률 (PFH)} & \textbf{요구 진단 커버리지} & \textbf{적용 사례} \\
\midrule
SIL 1 & $10^{-2}$ ~ $10^{-1}$ & $10^{-6}$ ~ $10^{-5}$/hr & $\geq$ 60\% & 일반 공정 경보, 온도 감시 \\
SIL 2 & $10^{-3}$ ~ $10^{-2}$ & $10^{-7}$ ~ $10^{-6}$/hr & $\geq$ 90\% & 압력 안전 시스템, 가스 탐지 \\
SIL 3 & $10^{-4}$ ~ $10^{-3}$ & $10^{-8}$ ~ $10^{-7}$/hr & $\geq$ 99\% & 비상 정지(ESD), 화재 방호 \\
SIL 4 & $10^{-5}$ ~ $10^{-4}$ & $10^{-9}$ ~ $10^{-8}$/hr & $\geq$ 99.9\% & 원자력, 항공 안전 시스템 \\
\bottomrule
\end{tabularx}
\end{table}

\paragraph{AI의 SIL 적용 한계}

IEC 61508은 본래 결정론적(Deterministic) 시스템을 전제로 설계되었다. AI/ML 모델은 본질적으로 확률적(Probabilistic)이며, 동일 입력에도 학습 조건에 따라 다른 출력을 생성할 수 있다. 이로 인해 AI 모델의 SIL 적용에는 근본적 한계가 존재한다.

\begin{table}[htbp]
\centering
\caption{AI 모델의 SIL 적용 한계와 대응 전략}
\label{tab:19_ai_sil_limits}
\begin{tabularx}{\textwidth}{p{3cm}Xp{4.5cm}}
\toprule
\textbf{한계 요인} & \textbf{설명} & \textbf{대응 전략} \\
\midrule
비결정론성 & 동일 입력에 다른 출력 가능 & 모델 고정(Frozen) 배포, 시드 제어 \\
블랙박스 특성 & 내부 로직 검증 불가능 & XAI 적용, 규칙 기반 래퍼(Wrapper) \\
데이터 의존성 & 학습 분포 외 입력 시 예측 불가 & ODD 제한, 분포 외 탐지 \\
드리프트\cite{paleyes2022challenges} & 시간 경과에 따른 성능 변화 & 연속 모니터링, 재검증 주기 설정 \\
테스트 완전성 & 모든 입력 조합 테스트 불가능 & 등가 클래스 테스트, 퍼즈 테스트 \\
소프트웨어 검증 & 기존 V\&V 방법론 한계 & 통계적 검증, 시뮬레이션 기반 테스트 \\
\bottomrule
\end{tabularx}
\end{table}

이러한 한계로 인해 현재 산업계의 합의는 \textbf{AI 단독으로 SIL 2 이상의 안전 기능을 수행해서는 안 된다}는 것이다. AI는 보조/감시 역할에 한정하고, 최종 안전 기능은 검증된 하드웨어 인터록(Hardware Interlock)이 담당하는 다층 방어 구조를 채택해야 한다.

\begin{lstlisting}[language=Python, caption={IEC 61508 AI Application Guide: SIL Compliance Assessment Framework}, label={lst:iec61508_guide}]
from dataclasses import dataclass, field
from typing import Dict, List, Optional, Tuple
from enum import Enum
from datetime import datetime

class SILLevel(Enum):
    SIL_1 = 1
    SIL_2 = 2
    SIL_3 = 3
    SIL_4 = 4

class AIInvolvement(Enum):
    DIRECT_CONTROL = "direct_control"
    ADVISORY = "advisory"
    MONITORING = "monitoring"
    ANALYTICS = "analytics"

class ComplianceStatus(Enum):
    COMPLIANT = "compliant"
    NON_COMPLIANT = "non_compliant"
    CONDITIONAL = "conditional"
    NOT_APPLICABLE = "not_applicable"

@dataclass
class SILAssessmentResult:
    function_name: str
    target_sil: SILLevel
    ai_involvement: AIInvolvement
    diagnostic_coverage: float
    hardware_interlock: bool
    compliance_status: ComplianceStatus
    conditions: List[str]
    risk_mitigations: List[str]

class IEC61508AIComplianceAssessor:
    """IEC 61508 compliance assessment for AI-integrated safety functions"""

    # Maximum SIL allowed per AI involvement type
    MAX_AI_SIL = {
        AIInvolvement.DIRECT_CONTROL: SILLevel.SIL_1,
        AIInvolvement.ADVISORY: SILLevel.SIL_2,
        AIInvolvement.MONITORING: SILLevel.SIL_3,
        AIInvolvement.ANALYTICS: SILLevel.SIL_4,  # No safety function
    }

    # Minimum diagnostic coverage per SIL
    MIN_DIAGNOSTIC_COVERAGE = {
        SILLevel.SIL_1: 0.60,
        SILLevel.SIL_2: 0.90,
        SILLevel.SIL_3: 0.99,
        SILLevel.SIL_4: 0.999,
    }

    def __init__(self):
        self.assessments: List[SILAssessmentResult] = []

    def assess_function(self, function_name: str, target_sil: SILLevel,
                        ai_involvement: AIInvolvement,
                        diagnostic_coverage: float,
                        hardware_interlock: bool,
                        odd_defined: bool,
                        drift_monitoring: bool) -> SILAssessmentResult:
        """Assess AI safety function against IEC 61508 requirements"""
        conditions = []
        mitigations = []
        max_sil = self.MAX_AI_SIL[ai_involvement]

        # Check SIL level compliance
        if target_sil.value > max_sil.value:
            conditions.append(
                f"Target SIL {target_sil.value} exceeds max "
                f"allowed SIL {max_sil.value} for {ai_involvement.value}"
            )
            if not hardware_interlock:
                conditions.append("Hardware interlock REQUIRED for higher SIL")
            else:
                mitigations.append("Hardware interlock provides independent safety layer")

        # Check diagnostic coverage
        min_dc = self.MIN_DIAGNOSTIC_COVERAGE[target_sil]
        if diagnostic_coverage < min_dc:
            conditions.append(
                f"Diagnostic coverage {diagnostic_coverage:.1%} "
                f"below required {min_dc:.1%} for SIL {target_sil.value}"
            )

        # Check ODD definition
        if not odd_defined:
            conditions.append("Operational Design Domain (ODD) not defined")
            mitigations.append("Define explicit ODD with boundary conditions")

        # Check drift monitoring
        if not drift_monitoring:
            conditions.append("No continuous drift monitoring system")
            mitigations.append("Implement PSI/KS-based drift detection")

        # Determine compliance
        critical_issues = [c for c in conditions if "REQUIRED" in c or "exceeds" in c]
        if not conditions:
            status = ComplianceStatus.COMPLIANT
        elif critical_issues and not hardware_interlock:
            status = ComplianceStatus.NON_COMPLIANT
        else:
            status = ComplianceStatus.CONDITIONAL

        result = SILAssessmentResult(
            function_name=function_name,
            target_sil=target_sil,
            ai_involvement=ai_involvement,
            diagnostic_coverage=diagnostic_coverage,
            hardware_interlock=hardware_interlock,
            compliance_status=status,
            conditions=conditions,
            risk_mitigations=mitigations
        )
        self.assessments.append(result)
        return result

    def generate_compliance_report(self) -> Dict:
        """Generate IEC 61508 AI compliance summary report"""
        total = len(self.assessments)
        compliant = sum(1 for a in self.assessments
                       if a.compliance_status == ComplianceStatus.COMPLIANT)
        conditional = sum(1 for a in self.assessments
                         if a.compliance_status == ComplianceStatus.CONDITIONAL)
        non_compliant = sum(1 for a in self.assessments
                           if a.compliance_status == ComplianceStatus.NON_COMPLIANT)

        return {
            "report_date": datetime.now().isoformat(),
            "total_functions": total,
            "compliant": compliant,
            "conditional": conditional,
            "non_compliant": non_compliant,
            "overall_status": "PASS" if non_compliant == 0 else "FAIL",
            "critical_findings": [
                {"function": a.function_name, "issues": a.conditions}
                for a in self.assessments
                if a.compliance_status == ComplianceStatus.NON_COMPLIANT
            ]
        }
\end{lstlisting}

\subsection{ISO 13849 PL 등급 상세}
\label{sec:19.2.2}

ISO 13849는 기계류의 안전 관련 제어 시스템 부품의 설계에 관한 표준이다 \cite{iso13849}. IEC 61508이 전체 안전 시스템의 라이프사이클을 다루는 반면, ISO 13849는 기계류 제어 시스템에 특화되어 있다.

\paragraph{PL(Performance Level) 등급}

ISO 13849는 안전 성능을 PL a~e의 5단계로 분류한다. PL은 위험 감소에 기여하는 제어 시스템의 능력을 나타내며, 위험성 평가를 통해 결정된다.

\begin{table}[htbp]
\centering
\caption{ISO 13849 PL 등급별 요구사항과 AI 적용 가능성}
\label{tab:19_pl}
\begin{tabularx}{\textwidth}{p{1cm}p{2.5cm}p{2cm}Xp{2.5cm}}
\toprule
\textbf{PL} & \textbf{위험측 고장 확률 (PFH$_\text{d}$)} & \textbf{카테고리} & \textbf{특성} & \textbf{AI 적용 가능성} \\
\midrule
PL a & $\geq 10^{-5}$ ~ $<10^{-4}$ & B, 1 & 기본 안전 원칙 적용 & AI 보조 가능 \\
PL b & $\geq 3 \times 10^{-6}$ ~ $<10^{-5}$ & 1, 2 & 검증된 안전 원칙 & AI 모니터링 가능 \\
PL c & $\geq 10^{-6}$ ~ $<3 \times 10^{-6}$ & 2, 3 & 단일 고장 검출 & AI 감시 (제한적) \\
PL d & $\geq 10^{-7}$ ~ $<10^{-6}$ & 3 & 단일 고장 허용 & AI 직접 제어 불가 \\
PL e & $\geq 10^{-8}$ ~ $<10^{-7}$ & 4 & 이중화, 모니터링 & AI 사용 제한 \\
\bottomrule
\end{tabularx}
\end{table}

AI를 기계류 안전 제어에 적용할 때는 PL c 이상에서 AI 단독 제어가 사실상 불가능하다. PL d/e에서는 하드웨어 이중화와 자기 진단 기능이 필수이며, 소프트웨어(AI 포함)는 보조적 역할에 한정된다.

\subsection{UL 4600 자율시스템 안전 평가 상세}
\label{sec:19.2.3}

UL 4600은 자율 시스템의 안전 평가를 위한 가이드라인으로, 자율 주행 차량에서 시작하여 산업용 자율 로봇, 드론, AMR(Autonomous Mobile Robot) 등으로 적용 범위가 확대되고 있다 \cite{ul4600}.

\paragraph{UL 4600 핵심 요구사항}

UL 4600은 전통적 SIL/PL 기반 표준과 달리, 자율 시스템의 특수성을 고려한 요구사항을 제시한다.

\begin{table}[htbp]
\centering
\caption{UL 4600 핵심 요구사항과 산업 AI 적용}
\label{tab:19_ul4600}
\begin{tabularx}{\textwidth}{p{3cm}Xp{4cm}}
\toprule
\textbf{요구사항 영역} & \textbf{내용} & \textbf{산업 AI 적용} \\
\midrule
안전 사례(Safety Case) & 구조적 논증을 통한 안전 입증 & 안전 사례 문서 작성 의무 \\
리스크 평가 & 잠재 리스크 식별 및 완화 & HAZOP/FMEA에 AI 리스크 추가 \\
ODD 정의 & 운영 설계 영역 명확화 & 작업 환경, 조건, 제한 명시 \\
데이터 무결성 & 학습/운영 데이터 품질 & 센서 데이터 검증, 라벨 품질 \\
소프트웨어 검증 & AI 모델 테스트 & 시뮬레이션 + 현장 테스트 \\
인간-기계 상호작용 & 작업자와의 안전한 협업 & HMI 설계, 비상 정지 \\
보안 & 사이버 보안 위협 대응 & 적대적 공격 방어, 통신 보안 \\
수명 주기 관리 & 배포 후 지속 관리 & 모니터링, 업데이트, 폐기 \\
\bottomrule
\end{tabularx}
\end{table}

\subsection{ISO/IEC 42001 AI 경영시스템}
\label{sec:19.2.4}

ISO/IEC 42001은 AI 경영시스템(AIMS, AI Management System) 인증 표준으로, 전사적 AI 관리 체계를 제공한다 \cite{iso42001}. 산업 현장에서는 기존 품질 관리 시스템(ISO 9001), 환경 관리 시스템(ISO 14001), 안전 관리 시스템(ISO 45001)과 통합하여 운영할 수 있다.

%----------------------------------------------------------
\section{구현 사례 연구}
\label{sec:19.3}

\subsection{예지보전(PdM) AI 거버넌스}
\label{sec:19.3.1}

예지보전은 장비 수명을 예측하여 다운타임을 최소화한다. 거버넌스는 \textbf{모델 드리프트} 탐지 체계와 AI 오판 시 수동 정비 절차의 일관성을 확보하는 데 집중해야 한다.

\paragraph{센서 데이터 관리 거버넌스}

예지보전 AI의 성능은 센서 데이터의 품질에 직접적으로 의존한다. 산업 현장의 센서는 극한 환경(고온, 진동, 분진)에서 작동하므로 데이터 품질 관리가 특히 중요하다.

\begin{table}[htbp]
\centering
\caption{예지보전 센서 데이터 관리 거버넌스 요건}
\label{tab:19_sensor_gov}
\begin{tabularx}{\textwidth}{p{2.5cm}Xp{4cm}}
\toprule
\textbf{관리 항목} & \textbf{내용} & \textbf{기준/절차} \\
\midrule
센서 캘리브레이션 & 센서 정확도 정기 교정 & 제조사 권장 주기 (최소 반기별) \\
결측 데이터 처리 & 센서 통신 장애 시 결측 관리 & 보간법 적용, 결측률 > 10\% 시 알림 \\
이상치 탐지 & 물리적으로 불가능한 값 제거 & 물리 법칙 기반 범위 설정 \\
샘플링 주기 & 적정 데이터 수집 빈도 & 장비 특성별 최소 샘플링 주기 설정 \\
데이터 동기화 & 다중 센서 시간 정합성 & NTP 동기화, 최대 허용 오차 < 100ms \\
라벨링 품질 & 고장/정상 라벨 정확성 & 정비 이력 대조, 전문가 검증 \\
\bottomrule
\end{tabularx}
\end{table}

\paragraph{드리프트 탐지와 알림 체계}

예지보전 AI에서 모델 드리프트는 장비 특성 변화(마모, 교체), 공정 조건 변경, 센서 열화 등 다양한 원인으로 발생한다. 드리프트를 조기에 탐지하고 적절히 대응하는 체계가 필수적이다.

\begin{lstlisting}[language=Python, caption={Predictive Maintenance AI Governance: Drift Detection and Alert System}, label={lst:pdm_drift}]
import numpy as np
import pandas as pd
from dataclasses import dataclass, field
from typing import Dict, List, Optional, Tuple
from datetime import datetime, timedelta
from enum import Enum
from collections import deque
from scipy import stats

class DriftSeverity(Enum):
    NORMAL = "normal"
    WARNING = "warning"
    CRITICAL = "critical"
    EMERGENCY = "emergency"

class AlertAction(Enum):
    LOG_ONLY = "log_only"
    NOTIFY_ENGINEER = "notify_engineer"
    SWITCH_TO_MANUAL = "switch_to_manual"
    HALT_PREDICTIONS = "halt_predictions"

@dataclass
class DriftAlert:
    alert_id: str
    equipment_id: str
    severity: DriftSeverity
    action: AlertAction
    drift_metric: str
    drift_value: float
    threshold: float
    timestamp: datetime
    description: str

class PdMAIDriftMonitor:
    """Predictive Maintenance AI drift detection and alert system"""

    SEVERITY_ACTIONS = {
        DriftSeverity.NORMAL: AlertAction.LOG_ONLY,
        DriftSeverity.WARNING: AlertAction.NOTIFY_ENGINEER,
        DriftSeverity.CRITICAL: AlertAction.SWITCH_TO_MANUAL,
        DriftSeverity.EMERGENCY: AlertAction.HALT_PREDICTIONS,
    }

    def __init__(self, equipment_id: str, window_size: int = 1000):
        self.equipment_id = equipment_id
        self.window_size = window_size
        self.baseline_distribution = None
        self.prediction_history = deque(maxlen=10000)
        self.alerts: List[DriftAlert] = []
        self.alert_counter = 0

        # Thresholds
        self.psi_warning = 0.10
        self.psi_critical = 0.25
        self.ks_warning = 0.05    # KS test p-value
        self.accuracy_drop_warning = 0.05
        self.accuracy_drop_critical = 0.10

    def set_baseline(self, baseline_predictions: np.ndarray,
                     baseline_features: np.ndarray):
        """Set baseline distribution from validation period"""
        self.baseline_distribution = {
            "predictions": baseline_predictions,
            "features": baseline_features,
            "mean": np.mean(baseline_predictions),
            "std": np.std(baseline_predictions),
            "set_date": datetime.now()
        }

    def monitor_batch(self, current_predictions: np.ndarray,
                      current_features: np.ndarray,
                      actual_outcomes: Optional[np.ndarray] = None) -> List[DriftAlert]:
        """Monitor a batch of predictions for drift"""
        alerts = []

        if self.baseline_distribution is None:
            return alerts

        # 1. PSI-based prediction drift detection
        psi = self._calculate_psi(
            self.baseline_distribution["predictions"],
            current_predictions
        )
        if psi >= self.psi_critical:
            alerts.append(self._create_alert(
                DriftSeverity.CRITICAL, "PSI", psi, self.psi_critical,
                f"Prediction distribution shift detected (PSI={psi:.4f})"
            ))
        elif psi >= self.psi_warning:
            alerts.append(self._create_alert(
                DriftSeverity.WARNING, "PSI", psi, self.psi_warning,
                f"Moderate prediction drift (PSI={psi:.4f})"
            ))

        # 2. KS test for feature drift
        for i in range(current_features.shape[1]):
            ks_stat, p_value = stats.ks_2samp(
                self.baseline_distribution["features"][:, i],
                current_features[:, i]
            )
            if p_value < self.ks_warning:
                alerts.append(self._create_alert(
                    DriftSeverity.WARNING, f"KS_Feature_{i}",
                    p_value, self.ks_warning,
                    f"Feature {i} distribution changed (KS p={p_value:.4f})"
                ))

        # 3. Accuracy-based drift (if actuals available)
        if actual_outcomes is not None:
            baseline_acc = 0.90  # Expected accuracy
            current_errors = np.abs(current_predictions - actual_outcomes)
            current_acc = 1.0 - np.mean(current_errors > 0.5)
            drop = baseline_acc - current_acc

            if drop >= self.accuracy_drop_critical:
                alerts.append(self._create_alert(
                    DriftSeverity.EMERGENCY, "Accuracy",
                    drop, self.accuracy_drop_critical,
                    f"Severe accuracy degradation ({drop:.2%} drop)"
                ))
            elif drop >= self.accuracy_drop_warning:
                alerts.append(self._create_alert(
                    DriftSeverity.WARNING, "Accuracy",
                    drop, self.accuracy_drop_warning,
                    f"Accuracy declining ({drop:.2%} drop)"
                ))

        self.alerts.extend(alerts)
        return alerts

    def _calculate_psi(self, expected, actual, buckets=10):
        """Calculate Population Stability Index"""
        exp_pct = np.histogram(expected, buckets)[0] / len(expected)
        act_pct = np.histogram(actual, buckets)[0] / len(actual)
        exp_pct = np.clip(exp_pct, 0.0001, 1)
        act_pct = np.clip(act_pct, 0.0001, 1)
        return float(np.sum((act_pct - exp_pct) * np.log(act_pct / exp_pct)))

    def _create_alert(self, severity, metric, value, threshold, desc):
        """Create a drift alert with appropriate action"""
        self.alert_counter += 1
        return DriftAlert(
            alert_id=f"DRIFT-{self.equipment_id}-{self.alert_counter:05d}",
            equipment_id=self.equipment_id,
            severity=severity,
            action=self.SEVERITY_ACTIONS[severity],
            drift_metric=metric,
            drift_value=value,
            threshold=threshold,
            timestamp=datetime.now(),
            description=desc
        )

    def get_status_summary(self) -> Dict:
        """Get current drift monitoring status"""
        recent_alerts = [a for a in self.alerts
                        if (datetime.now() - a.timestamp).days < 7]
        max_severity = max(
            [a.severity for a in recent_alerts],
            key=lambda s: list(DriftSeverity).index(s),
            default=DriftSeverity.NORMAL
        )
        return {
            "equipment_id": self.equipment_id,
            "status": max_severity.value,
            "recent_alerts_7d": len(recent_alerts),
            "total_alerts": len(self.alerts),
            "recommended_action": self.SEVERITY_ACTIONS[max_severity].value,
            "baseline_age_days": (datetime.now() - self.baseline_distribution["set_date"]).days
                                 if self.baseline_distribution else None
        }
\end{lstlisting}

\subsection{품질 검사 AI 거버넌스}
\label{sec:19.3.2}

AI 기반 품질 검사(Vision AI)는 제조업에서 가장 빠르게 확산되는 AI 응용 영역이다. 제품의 외관 결함, 치수 이상, 조립 오류 등을 자동으로 탐지하며, 인간 검사원의 피로와 주관성 문제를 해결한다.

\paragraph{품질 검사 AI의 핵심 거버넌스 과제}

\begin{enumerate}
 \item \textbf{과검(Over-kill) vs 미검(Escape)의 균형}: 과검은 양품을 불량으로 판정하여 수율을 낮추고, 미검은 불량품을 시장에 출하하여 품질 문제와 리콜 비용을 야기한다. 두 오류의 비용 구조가 다르므로, 단순 정확도 최적화가 아닌 비용 최적화가 필요하다.
 \item \textbf{검사 기준의 일관성}: AI 모델이 업데이트될 때마다 검사 기준이 미묘하게 변할 수 있어, 품질 이력의 연속성이 훼손될 수 있다.
 \item \textbf{검사 추적성(Traceability)}: 특정 제품이 왜 합격/불합격 판정을 받았는지 사후 추적이 가능해야 한다. 고객 클레임, 리콜 대응, 규제 감사에 필수적이다.
\end{enumerate}

\begin{table}[htbp]
\centering
\caption{과검과 미검의 비용 구조}
\label{tab:19_8}
\begin{tabularx}{\textwidth}{p{3cm}p{3cm}XX}
\toprule
\textbf{오류 유형} & \textbf{정의} & \textbf{직접 비용} & \textbf{간접 비용} \\
\midrule
\textbf{과검 (Over-kill)} & 양품을 불량 판정 & 재작업/Retire 비용, 수율 저하 & 생산성 하락, 납기 지연 \\
\textbf{미검 (Escape)} & 불량을 양품 판정 & 고객 클레임, 리콜 비용 & 브랜드 훼손, 품질 인증 위험 \\
\bottomrule
\end{tabularx}
\end{table}

일반적으로 \textbf{미검 비용이 과검 비용의 10~100배}에 달한다. 자동차 부품의 경우 단일 불량이 대규모 리콜로 이어질 수 있어 미검 비용이 극도로 높다. 반면 반도체나 디스플레이 산업에서는 수율이 수익성에 직결되어 과검 비용도 무시할 수 없다.

\paragraph{검사 기준 일관성 관리}

AI 모델이 업데이트(재학습)될 때 검사 기준이 변하는 문제를 관리하기 위한 체계가 필요하다. 모델 업데이트 전후의 검사 결과를 동일 샘플로 비교하는 '골든 샘플(Golden Sample) 테스트'가 핵심이다.

\begin{table}[htbp]
\centering
\caption{품질 검사 AI 모델 업데이트 관리 체계}
\label{tab:19_model_update}
\begin{tabularx}{\textwidth}{p{3cm}Xp{3.5cm}}
\toprule
\textbf{관리 항목} & \textbf{내용} & \textbf{기준} \\
\midrule
골든 샘플 테스트 & 표준 불량/양품 세트로 신구 모델 비교 & 불일치율 < 2\% \\
판정 일관성 & 동일 제품 반복 판정 시 일치율 & 반복 일치율 > 99\% \\
경계 사례 관리 & 경계 영역(불량/양품 구분 어려움) 별도 관리 & 인간 검토 비율 설정 \\
이력 비교 & 모델 변경\cite{kreuzberger2023machine} 전후 과검/미검률 비교 & 미검률 증가 불허 \\
단계적 적용 & 신모델 병렬 운영 후 교체 & 최소 2주 병렬 운영 \\
\bottomrule
\end{tabularx}
\end{table}

\paragraph{검사 이력 추적 체계}

품질 검사 AI의 모든 판정은 추적 가능(Traceable)해야 한다. 이는 고객 클레임 대응, 리콜 범위 결정, 규제 감사에 필수적이다.

\begin{lstlisting}[language=Python, caption={Quality Inspection AI Governance: Cost Optimization and Traceability Management}, label={lst:quality_gov_full}]
from dataclasses import dataclass, field
from typing import Dict, List, Optional, Tuple, Any
from enum import Enum
from datetime import datetime
import numpy as np
from scipy.optimize import minimize_scalar

class DefectClass(Enum):
 CRITICAL = "Critical" # Impact on safety/function
 MAJOR = "Major" # Potential performance degradation
 MINOR = "Minor" # Appearance defect
 NONE = "None"

class InspectionDecision(Enum):
 PASS = "Pass"
 FAIL = "Fail"
 REINSPECT = "Re-inspect"
 HUMAN_REVIEW = "Worker Review"

@dataclass
class QualityCostStructure:
 """Quality Cost Structure"""
 overkill_cost_per_unit: float # Unit rework/scrap cost
 escape_cost_per_unit: float # Unit claim/recall cost
 escape_recall_probability: float # Recall probability
 recall_cost_multiplier: float # Cost multiplier on recall

class VisionAIQualityGovernance:
 """Vision AI Quality Inspection Governance Manager"""
 def __init__(self, cost_structure: QualityCostStructure, threshold=0.5):
 self.cost_structure = cost_structure
 self.current_threshold = threshold
 self.inspection_history = []

 def optimize_threshold(self, predictions, true_labels):
 cs = self.cost_structure
 def calculate_total_cost(t):
 y_pred = predictions >= t
 fp = np.sum((y_pred == 1) & (true_labels == 0)) # Overkill
 fn = np.sum((y_pred == 0) & (true_labels == 1)) # Escape

 overkill_cost = fp * cs.overkill_cost_per_unit
 escape_cost = fn * cs.escape_cost_per_unit * (1 + cs.escape_recall_probability * cs.recall_cost_multiplier)
 return overkill_cost + escape_cost

 result = minimize_scalar(calculate_total_cost, bounds=(0.1, 0.9), method='bounded')
 return result.x

 def perform_inspection(self, product_id, prob, defect_class, confidence):
 # Decision logic: based on threshold and confidence
 if prob >= self.current_threshold:
 decision = InspectionDecision.FAIL if confidence > 0.7 else InspectionDecision.HUMAN_REVIEW
 else:
 decision = InspectionDecision.PASS

 result = {"product_id": product_id, "decision": decision, "threshold": self.current_threshold}
 self.inspection_history.append(result)
 return decision
\end{lstlisting}

\begin{lstlisting}[language=Python, caption={Quality Inspection Traceability and History Tracking System}, label={lst:quality_trace}]
import hashlib
import json
from dataclasses import dataclass, field
from typing import Dict, List, Optional
from datetime import datetime

@dataclass
class InspectionRecord:
    record_id: str
    product_id: str
    lot_id: str
    line_id: str
    model_version: str
    timestamp: datetime
    raw_probability: float
    threshold_used: float
    decision: str
    defect_class: str
    confidence: float
    image_hash: str
    explanation: Dict
    operator_override: Optional[str] = None
    override_reason: Optional[str] = None

class QualityTraceabilitySystem:
    """Full traceability system for quality inspection AI"""

    def __init__(self, retention_years: int = 10):
        self.records: Dict[str, InspectionRecord] = {}
        self.retention_years = retention_years
        self.golden_samples: Dict[str, Dict] = {}

    def record_inspection(self, product_id: str, lot_id: str,
                          line_id: str, model_version: str,
                          probability: float, threshold: float,
                          decision: str, defect_class: str,
                          confidence: float, image_bytes: bytes,
                          explanation: Dict) -> InspectionRecord:
        """Record complete inspection with full traceability"""
        image_hash = hashlib.sha256(image_bytes).hexdigest()
        record_id = f"INS-{datetime.now().strftime('%Y%m%d%H%M%S')}-{product_id}"

        record = InspectionRecord(
            record_id=record_id,
            product_id=product_id,
            lot_id=lot_id,
            line_id=line_id,
            model_version=model_version,
            timestamp=datetime.now(),
            raw_probability=probability,
            threshold_used=threshold,
            decision=decision,
            defect_class=defect_class,
            confidence=confidence,
            image_hash=image_hash,
            explanation=explanation
        )
        self.records[record_id] = record
        return record

    def record_operator_override(self, record_id: str,
                                  new_decision: str, reason: str):
        """Record when operator overrides AI decision"""
        if record_id in self.records:
            self.records[record_id].operator_override = new_decision
            self.records[record_id].override_reason = reason

    def trace_by_product(self, product_id: str) -> List[InspectionRecord]:
        """Retrieve all inspection records for a product"""
        return [r for r in self.records.values() if r.product_id == product_id]

    def trace_by_lot(self, lot_id: str) -> List[InspectionRecord]:
        """Retrieve all records for a lot (recall scope determination)"""
        return [r for r in self.records.values() if r.lot_id == lot_id]

    def golden_sample_test(self, model_version: str,
                            samples: Dict[str, float]) -> Dict:
        """Run golden sample test for model update validation"""
        if not self.golden_samples:
            self.golden_samples = samples
            return {"status": "baseline_set", "model": model_version}

        mismatches = 0
        total = len(self.golden_samples)
        for sample_id, expected in self.golden_samples.items():
            actual = samples.get(sample_id, None)
            if actual is None or abs(actual - expected) > 0.1:
                mismatches += 1

        mismatch_rate = mismatches / total if total > 0 else 0
        return {
            "model_version": model_version,
            "total_samples": total,
            "mismatches": mismatches,
            "mismatch_rate": mismatch_rate,
            "status": "PASS" if mismatch_rate < 0.02 else "FAIL",
            "test_date": datetime.now().isoformat()
        }

    def generate_audit_report(self, start_date: datetime,
                               end_date: datetime) -> Dict:
        """Generate inspection audit report for regulatory compliance"""
        period_records = [
            r for r in self.records.values()
            if start_date <= r.timestamp <= end_date
        ]
        total = len(period_records)
        overrides = sum(1 for r in period_records if r.operator_override)

        return {
            "period": f"{start_date.date()} ~ {end_date.date()}",
            "total_inspections": total,
            "pass_count": sum(1 for r in period_records if r.decision == "Pass"),
            "fail_count": sum(1 for r in period_records if r.decision == "Fail"),
            "human_review_count": sum(1 for r in period_records if r.decision == "Worker Review"),
            "operator_overrides": overrides,
            "override_rate": overrides / total if total > 0 else 0,
            "model_versions_used": list(set(r.model_version for r in period_records))
        }
\end{lstlisting}

\subsection{자율로봇/AMR 거버넌스}
\label{sec:19.3.3}

자율 이동 로봇(AMR, Autonomous Mobile Robot)과 협동 로봇(Cobot)은 스마트 팩토리의 핵심 요소로 급속히 확산되고 있다. 이들 로봇은 작업자와 동일 공간에서 작동하므로, 안전 거버넌스가 특히 중요하다.

\paragraph{안전 영역(Safety Zone) 관리}

자율로봇의 안전 거버넌스에서 가장 기본적인 개념은 안전 영역의 정의이다. 로봇의 속도와 행동을 영역별로 제한하여 작업자 안전을 확보한다.

\begin{table}[htbp]
\centering
\caption{자율로봇 안전 영역 체계}
\label{tab:19_safety_zones}
\begin{tabularx}{\textwidth}{p{2.5cm}p{2cm}Xp{3.5cm}}
\toprule
\textbf{영역 유형} & \textbf{작업자 존재} & \textbf{로봇 행동 제한} & \textbf{AI 관련 요건} \\
\midrule
자유 이동 영역 & 없음 & 최대 속도 운행 & 기본 경로 계획 AI \\
공유 영역 & 가능 & 속도 제한 (1.5m/s) & 작업자 인식 AI 필수 \\
협업 영역 & 상시 & 저속 + 힘 제한 & 실시간 거리 모니터링 \\
금지 영역 & 해당 없음 & 진입 금지 & 가상 울타리 AI \\
비상 영역 & 긴급 & 즉시 정지 & 비상 정지 우선 (하드웨어) \\
\bottomrule
\end{tabularx}
\end{table}

\paragraph{인간-로봇 협업(HRC) 거버넌스}

인간-로봇 협업(Human-Robot Collaboration, HRC) 환경에서는 ISO 10218(산업용 로봇 안전)과 ISO/TS 15066(협동 로봇 안전)을 준수해야 한다. AI가 HRC에 적용될 때의 거버넌스 요건은 다음과 같다:

\begin{itemize}
    \item \textbf{작업자 인식}: AI 기반 인체 탐지 시스템의 정확도는 99.9\% 이상이어야 하며, 미탐지 시 즉시 안전 정지가 작동해야 한다.
    \item \textbf{의도 예측}: 작업자의 이동 방향과 의도를 예측하여 선제적으로 회피하는 AI는 보조 기능으로만 사용하고, 최종 안전은 물리적 센서(라이다, 안전 스캐너)가 담당한다.
    \item \textbf{힘/속도 제한}: 협업 모드에서 로봇의 최대 접촉력과 속도는 ISO/TS 15066에 따라 인체 부위별로 제한되며, AI가 이 한계를 초과하는 명령을 생성할 수 없도록 물리적 리미터를 설정한다.
    \item \textbf{비상 정지}: AI 시스템 장애 시에도 비상 정지는 하드와이어드(Hardwired) 회로로 독립 작동해야 한다.
\end{itemize}

\paragraph{자율로봇의 ODD(Operational Design Domain) 관리}

UL 4600에 따라 자율로봇의 운영 설계 영역(ODD)을 명확히 정의하고 관리하는 것은 거버넌스의 핵심이다.

\begin{table}[htbp]
\centering
\caption{자율로봇 ODD 정의 항목}
\label{tab:19_odd}
\begin{tabularx}{\textwidth}{p{3cm}Xp{4cm}}
\toprule
\textbf{ODD 항목} & \textbf{정의 내용} & \textbf{ODD 이탈 시 대응} \\
\midrule
작업 환경 & 실내/실외, 바닥 유형, 경사도 & 안전 정지 후 대기 \\
조명 조건 & 최소/최대 조도, 역광 & 저속 운행 또는 정지 \\
장애물 유형 & 정적/동적, 크기 범위 & 미인식 장애물 시 정지 \\
작업자 밀도 & 최대 허용 인원 수 & 밀도 초과 시 대기 \\
통신 환경 & Wi-Fi/5G 신호 강도 & 통신 두절 시 안전 정지 \\
적재 조건 & 최대 중량, 무게 중심 & 초과 시 운반 거부 \\
온도/습도 & 허용 작동 범위 & 범위 이탈 시 경고/정지 \\
\bottomrule
\end{tabularx}
\end{table}

\subsection{디지털 트윈 거버넌스}
\label{sec:19.3.3b}

디지털 트윈(Digital Twin)은 물리적 자산 또는 프로세스의 가상 복제본으로, AI와 결합하여 시뮬레이션, 예측, 최적화에 활용된다. 디지털 트윈 거버넌스의 핵심 과제는 \textbf{가상-물리 일치성(Virtual-Physical Consistency)}의 유지이다.

\paragraph{가상-물리 일치성 검증}

디지털 트윈이 실제 물리 시스템을 정확하게 반영하지 못하면, 이에 기반한 AI 의사결정은 위험할 수 있다. 일치성 검증은 지속적으로 수행되어야 한다.

\begin{table}[htbp]
\centering
\caption{디지털 트윈 가상-물리 일치성 검증 체계}
\label{tab:19_dt_validation}
\begin{tabularx}{\textwidth}{p{2.5cm}Xp{3cm}p{2cm}}
\toprule
\textbf{검증 항목} & \textbf{내용} & \textbf{허용 오차} & \textbf{점검 주기} \\
\midrule
정적 파라미터 & 장비 사양, 배치, 연결 정보 & 0\% (완전 일치) & 변경 시 즉시 \\
동적 상태 & 온도, 압력, 속도 등 실시간 값 & $\pm$ 5\% & 실시간 \\
물리 모델 & 열역학, 유체 역학 시뮬레이션 & $\pm$ 10\% & 월간 \\
AI 예측 & 예지보전, 품질 예측 결과 & RMSE 기준 & 주간 \\
시나리오 응답 & What-if 시뮬레이션 정확도 & $\pm$ 15\% & 분기별 \\
\bottomrule
\end{tabularx}
\end{table}

\begin{lstlisting}[language=Python, caption={Digital Twin Governance: Virtual-Physical Consistency Validation}, label={lst:dt_validation}]
import numpy as np
from dataclasses import dataclass, field
from typing import Dict, List, Optional
from datetime import datetime
from enum import Enum

class ConsistencyStatus(Enum):
    SYNCHRONIZED = "synchronized"
    DRIFTING = "drifting"
    DESYNCHRONIZED = "desynchronized"
    CRITICAL = "critical"

@dataclass
class ConsistencyCheckResult:
    parameter: str
    physical_value: float
    virtual_value: float
    deviation: float
    tolerance: float
    status: ConsistencyStatus
    timestamp: datetime

class DigitalTwinGovernance:
    """Governance framework for digital twin virtual-physical consistency"""

    def __init__(self, asset_id: str):
        self.asset_id = asset_id
        self.check_history: List[ConsistencyCheckResult] = []
        self.tolerances: Dict[str, float] = {}
        self.calibration_log: List[Dict] = []

    def set_tolerances(self, tolerances: Dict[str, float]):
        """Set acceptable deviation tolerances per parameter"""
        self.tolerances = tolerances

    def check_consistency(self, physical_data: Dict[str, float],
                          virtual_data: Dict[str, float]) -> List[ConsistencyCheckResult]:
        """Check consistency between physical and virtual states"""
        results = []

        for param in physical_data:
            if param not in virtual_data:
                continue

            phys = physical_data[param]
            virt = virtual_data[param]
            deviation = abs(phys - virt) / (abs(phys) + 1e-10)
            tolerance = self.tolerances.get(param, 0.05)

            if deviation < tolerance:
                status = ConsistencyStatus.SYNCHRONIZED
            elif deviation < tolerance * 2:
                status = ConsistencyStatus.DRIFTING
            elif deviation < tolerance * 4:
                status = ConsistencyStatus.DESYNCHRONIZED
            else:
                status = ConsistencyStatus.CRITICAL

            result = ConsistencyCheckResult(
                parameter=param,
                physical_value=phys,
                virtual_value=virt,
                deviation=deviation,
                tolerance=tolerance,
                status=status,
                timestamp=datetime.now()
            )
            results.append(result)

        self.check_history.extend(results)
        return results

    def requires_recalibration(self) -> bool:
        """Determine if digital twin needs recalibration"""
        recent = [c for c in self.check_history
                 if (datetime.now() - c.timestamp).hours < 24]

        desync_count = sum(1 for c in recent
                          if c.status in [ConsistencyStatus.DESYNCHRONIZED,
                                          ConsistencyStatus.CRITICAL])
        return desync_count > len(recent) * 0.1 if recent else False

    def log_calibration(self, method: str, params_updated: List[str]):
        """Log recalibration event"""
        self.calibration_log.append({
            "timestamp": datetime.now().isoformat(),
            "method": method,
            "params_updated": params_updated,
            "pre_calibration_status": self._get_overall_status()
        })

    def _get_overall_status(self) -> str:
        """Get overall consistency status"""
        if not self.check_history:
            return "unknown"
        recent = self.check_history[-10:]
        statuses = [c.status for c in recent]
        if ConsistencyStatus.CRITICAL in statuses:
            return "critical"
        if ConsistencyStatus.DESYNCHRONIZED in statuses:
            return "desynchronized"
        return "healthy"

    def generate_governance_report(self) -> Dict:
        """Generate digital twin governance report"""
        total_checks = len(self.check_history)
        sync_rate = sum(1 for c in self.check_history
                       if c.status == ConsistencyStatus.SYNCHRONIZED) / total_checks \
                   if total_checks > 0 else 0

        return {
            "asset_id": self.asset_id,
            "report_date": datetime.now().isoformat(),
            "total_consistency_checks": total_checks,
            "synchronization_rate": sync_rate,
            "calibrations_performed": len(self.calibration_log),
            "current_status": self._get_overall_status(),
            "needs_recalibration": self.requires_recalibration()
        }
\end{lstlisting}

\subsection{해양/조선 AI 거버넌스 (HD현대Heavy Industries 사례)}
\label{sec:19.3.4}

\begin{practicaltool}{스마트십야드 AI 거버넌스: HD현대Heavy Industries Case}
조선소의 극한 환경(진동, 분진) 하에서 운용되는 AI 시스템의 신뢰성을 확보한 사례이다.\\
\textbf{용접 품질 AI}: 실시간 용접 파라미터 분석을 통해 결함을 예측하며, 주의 이상 상황 시 용접사 HMI에 즉시 경보를 송출하되, 최종 결정은 작업자가 수행한다.\\
\textbf{크레인 안전}: 대형 블록 인양 시 무게 중심과 풍속을 분석하여 안정성 여유(Stability Margin)를 계산한다. AI는 관제사에게 정지 권고를 내리며, 물리적 긴급 정지는 SIL 2 인증 장치가 담당한다. 특히 \textbf{뉴로심볼릭(Neuro-symbolic) AI} 접근법을 적용하여, 딥러닝의 시각 지능과 '블록 하중 제한 물리 법칙'이라는 심볼릭 규칙을 결합함으로써 의사결정의 신뢰도를 극대화하였다(20.1.3절 참조).\\
\textbf{Data 거버넌스}: 현장 작업 로그와 AI 판정 이력을 10년 이상 보관하여 품질 사고 시 리니지(계보) 추적이 가능하도록 설계되었다.
\end{practicaltool}

\section{추가 산업별 사례 및 체크리스트}
\label{sec:19.4}

\subsection{H제철: 고로 제어 AI의 물리적 Guardrail}
H제철은 쇳물 온도를 예측하고 연료 투입량을 자동 제어하는 AI 모델을 도입하면서, 센서 오작동으로 인한 폭발 위험을 방지하기 위해 '물리적 한계(Rule-based)'를 적용한 이중 안전 장치를 마련하였다.

\begin{table}[htbp]
\centering
\caption{H제철 고로 AI 물리적 가드레일 설정 예시}
\label{tab:19_guardrail}
\begin{tabularx}{\textwidth}{p{3cm}p{2.5cm}p{2.5cm}X}
\toprule
\textbf{제어 변수} & \textbf{AI 허용 범위} & \textbf{물리적 한계} & \textbf{한계 초과 시 조치} \\
\midrule
쇳물 온도 & 1,480--1,520\textdegree C & 1,450--1,550\textdegree C & 자동 경보 + 수동 전환 \\
풍량 & $\pm$10\% 조정 & $\pm$20\% 이내 & 하드웨어 리미터 작동 \\
코크스 투입량 & $\pm$5\% 조정 & $\pm$15\% 이내 & 투입 밸브 잠금 \\
노내 압력 & $\pm$0.05 atm 조정 & $\pm$0.15 atm 이내 & 비상 감압 밸브 \\
\bottomrule
\end{tabularx}
\end{table}

\subsection{S전자: 반도체 품질 검사의 드리프트 관리}
반도체 웨이퍼 불량 검출 시 발생하는 데이터 드리프트 문제를 해결하기 위해 MLOps 파이프라인에 상시 성능 모니터링 체계를 구축하고, 성능 저하 시 즉시 육안 검사 모드로 전환하는 프로세스를 표준화하였다.

\begin{table}[htbp]
\centering
\caption{S전자 반도체 품질 AI 드리프트 대응 체계}
\label{tab:19_samsung_drift}
\begin{tabularx}{\textwidth}{p{2.5cm}p{2cm}Xp{3cm}}
\toprule
\textbf{드리프트 원인} & \textbf{탐지 방법} & \textbf{발현 사례} & \textbf{대응 절차} \\
\midrule
공정 변경 & PSI 모니터링 & 새 레시피 적용 시 패턴 변화 & 모델 재학습 트리거 \\
장비 교체 & 특성 분포 비교 & 새 노광기 도입 시 이미지 특성 변화 & 전이 학습 적용 \\
소재 변경 & 이상치 탐지 & 새 웨이퍼 소재의 반사율 차이 & 골든 샘플 재설정 \\
센서 열화 & 보정값 추이 & 카메라 렌즈 오염, 조명 노화 & 캘리브레이션 수행 \\
계절 변화 & 장기 추세 분석 & 온습도 변화에 따른 검사 편차 & 환경 보정 모델 적용 \\
\bottomrule
\end{tabularx}
\end{table}

%----------------------------------------------------------
\section{산업 AI 안전성 평가 종합 체크리스트}
\label{sec:19.5}

\begin{practicaltool}{산업 AI 안전성 평가 종합 체크리스트 (확장)}
\textbf{1. 기능 안전 (Functional Safety)}
\begin{itemize}
    \item AI 출력이 물리적 리미터를 벗어나지 않는가? \checkmark / \texttimes
    \item 적용 SIL/PL 등급이 적절히 평가되었는가? \checkmark / \texttimes
    \item AI 단독으로 SIL 2 이상 안전 기능을 수행하지 않는가? \checkmark / \texttimes
    \item 하드웨어 인터록(비상 정지 등)이 독립 작동하는가? \checkmark / \texttimes
    \item 다층 방어(Defense in Depth) 구조가 적용되었는가? \checkmark / \texttimes
\end{itemize}

\textbf{2. 페일세이프(Fail-Safe) 설계}
\begin{itemize}
    \item 통신 두절 시 즉시 안전 모드로 전환되는가? \checkmark / \texttimes
    \item AI 모델 오류 시 수동 운전 전환 절차가 있는가? \checkmark / \texttimes
    \item 전원 상실 시 안전 상태가 보장되는가? \checkmark / \texttimes
    \item 비상 정지 기능이 하드와이어드로 구현되었는가? \checkmark / \texttimes
\end{itemize}

\textbf{3. ODD(Operational Design Domain) 관리}
\begin{itemize}
    \item 운영 설계 영역이 명확히 정의되었는가? \checkmark / \texttimes
    \item ODD 이탈 탐지 메커니즘이 있는가? \checkmark / \texttimes
    \item ODD 이탈 시 안전한 대체 행동이 정의되었는가? \checkmark / \texttimes
    \item ODD 정의가 정기적으로 검토/갱신되는가? \checkmark / \texttimes
\end{itemize}

\textbf{4. 데이터 강건성}
\begin{itemize}
    \item 센서 노이즈와 결측치가 충분히 반영되었는가? \checkmark / \texttimes
    \item 센서 캘리브레이션이 정기적으로 수행되는가? \checkmark / \texttimes
    \item 데이터 드리프트 모니터링 체계가 구축되었는가? \checkmark / \texttimes
    \item 적대적 입력에 대한 방어가 고려되었는가? \checkmark / \texttimes
\end{itemize}

\textbf{5. 인간-기계 상호작용(HMI)}
\begin{itemize}
    \item 작업자가 AI의 결정을 직관적으로 이해하고 제어할 수 있는가? \checkmark / \texttimes
    \item AI 권고와 작업자 최종 결정 간 역할이 명확한가? \checkmark / \texttimes
    \item 작업자 오버라이드(Override) 기능이 있는가? \checkmark / \texttimes
    \item AI 상태 표시(정상/경고/위험)가 명확한가? \checkmark / \texttimes
\end{itemize}

\textbf{6. OT-IT 보안}
\begin{itemize}
    \item 퍼듀 모델 기반 계층 분리가 적용되었는가? \checkmark / \texttimes
    \item 산업 DMZ가 구축되었는가? \checkmark / \texttimes
    \item 데이터 다이오드 또는 단방향 게이트웨이가 적용되었는가? \checkmark / \texttimes
    \item IEC 62443 보안 수준이 적정하게 설정되었는가? \checkmark / \texttimes
    \item AI 모델 업데이트 경로가 보안 통제되는가? \checkmark / \texttimes
\end{itemize}

\textbf{7. 모니터링 및 감사}
\begin{itemize}
    \item 성능 모니터링이 실시간으로 수행되는가? \checkmark / \texttimes
    \item 모든 AI 판정 이력이 추적 가능한가? \checkmark / \texttimes
    \item 정기 감사 주기가 설정되어 있는가? \checkmark / \texttimes
    \item 감사 보고서가 문서화되어 보관되는가? \checkmark / \texttimes
\end{itemize}

\textbf{8. 자율로봇/AMR 추가 점검}
\begin{itemize}
    \item 안전 영역(Safety Zone)이 정의되고 시행되는가? \checkmark / \texttimes
    \item 작업자 인식 AI 정확도가 99.9\% 이상인가? \checkmark / \texttimes
    \item 협업 모드에서 힘/속도 제한이 ISO/TS 15066을 준수하는가? \checkmark / \texttimes
    \item ODD 이탈 시 안전 정지가 자동으로 작동하는가? \checkmark / \texttimes
\end{itemize}

\textbf{9. 디지털 트윈 추가 점검}
\begin{itemize}
    \item 가상-물리 일치성 검증이 정기적으로 수행되는가? \checkmark / \texttimes
    \item 일치성 이탈 시 재캘리브레이션 절차가 있는가? \checkmark / \texttimes
    \item 디지털 트윈 기반 AI 의사결정이 물리 시스템에 직접 반영되기 전 인간 검토가 있는가? \checkmark / \texttimes
\end{itemize}
\end{practicaltool}

\subsection{에너지/발전 AI 거버넌스}
\label{sec:19.4.3}

에너지 및 발전 분야의 AI는 전력망 최적화, 재생에너지 출력 예측, 발전 설비 예지보전 등에 활용된다. 이 분야는 국가 기간시설이라는 특수성으로 인해 더욱 엄격한 거버넌스가 요구된다.

\begin{table}[htbp]
\centering
\caption{에너지/발전 AI 유형별 거버넌스 요건}
\label{tab:19_energy}
\begin{tabularx}{\textwidth}{p{3cm}p{2cm}Xp{3cm}}
\toprule
\textbf{AI 유형} & \textbf{안전 등급} & \textbf{핵심 리스크} & \textbf{거버넌스 요건} \\
\midrule
전력 수요 예측 & 중간 & 과소 예측 시 정전 위험 & 예측 오차 한도, 수동 보정 체계 \\
재생에너지 출력 예측 & 중간 & 예측 오류 시 전력망 불안정 & 앙상블 모델, 불확실성 정량화 \\
발전 설비 예지보전 & 높음 & 미탐지 고장 시 대정전 & SIL 2 이상 독립 안전 장치 \\
전력망 자동 제어 & 매우 높음 & 오제어 시 광역 정전 & AI 보조만 허용, 인간 최종 제어 \\
에너지 거래 최적화 & 중간 & 시장 교란, 과도 투기 & 거래 한도 설정, 이상 거래 탐지 \\
\bottomrule
\end{tabularx}
\end{table}

\paragraph{전력망 AI의 다층 방어 구조}

전력망 AI는 국가 기간시설의 안정성과 직결되므로, 다층 방어(Defense in Depth) 구조가 특히 중요하다. AI의 제어 명령은 다음의 검증 계층을 통과해야 한다:

\begin{enumerate}
    \item \textbf{AI 모델 출력}: 최적화 알고리즘이 제어 명령을 생성
    \item \textbf{물리 법칙 검증}: 생성된 명령이 전력 공학 법칙(전압/주파수 범위, 송전 용량 등)을 위반하지 않는지 규칙 기반 검증
    \item \textbf{안전 여유 확인}: 명령 수행 후 시스템 안전 여유(Safety Margin)가 기준값 이상인지 확인
    \item \textbf{운영자 승인}: 대규모 제어 변경은 운영자의 최종 승인 후 실행
    \item \textbf{하드웨어 보호}: 과전압/과전류 보호 계전기가 최종 안전 기능 수행
\end{enumerate}

\subsection{화학/석유화학 AI 거버넌스}
\label{sec:19.4.4}

화학 및 석유화학 공정의 AI는 공정 최적화, 이상 탐지, 안전 예측에 활용된다. 이 분야는 폭발, 화재, 유독가스 누출 등의 중대 사고 위험이 존재하므로, 공정안전관리(PSM, Process Safety Management)와의 통합이 핵심이다.

\begin{table}[htbp]
\centering
\caption{화학/석유화학 AI 거버넌스 핵심 요건}
\label{tab:19_chemical}
\begin{tabularx}{\textwidth}{p{3cm}Xp{4cm}}
\toprule
\textbf{거버넌스 항목} & \textbf{내용} & \textbf{관련 규정/표준} \\
\midrule
공정 변경 관리(MOC) & AI 모델 변경을 공정 변경으로 취급 & PSM 규정, OSHA 1910.119 \\
안전 계장 시스템(SIS) & AI는 SIS와 독립, SIS는 SIL 인증 & IEC 61511 \\
HAZOP 연계 & AI 도입 시 HAZOP 재수행 & IEC 61882 \\
비상 대응 & AI 장애 시 비상 운전 절차 & 비상대응계획(ERP) \\
환경 모니터링 & 배출 예측 AI의 정확도 관리 & 대기환경보전법 \\
\bottomrule
\end{tabularx}
\end{table}

중요한 점은 화학 공정에서 AI 모델의 변경(재학습, 파라미터 수정)은 \textbf{공정 변경 관리(Management of Change, MOC)} 절차를 따라야 한다는 것이다. 이는 AI 모델 변경이 공정 조건 변경과 동일한 위험을 초래할 수 있기 때문이다. MOC 절차는 변경 전 위험성 평가, 관련 부서 협의, 운영 절차서 갱신, 작업자 교육, 변경 후 검증의 단계를 포함한다.

\subsection{산업 AI 거버넌스 조직 구조}
\label{sec:19.4.5}

산업 AI 거버넌스를 효과적으로 운영하기 위해서는 기존 안전 관리 조직과 AI 거버넌스 조직을 통합하는 구조가 필요하다.

\begin{table}[htbp]
\centering
\caption{산업 AI 거버넌스 조직 역할 분담}
\label{tab:19_org}
\begin{tabularx}{\textwidth}{p{3cm}p{3cm}X}
\toprule
\textbf{조직} & \textbf{기존 역할} & \textbf{AI 거버넌스 추가 역할} \\
\midrule
안전 관리실 & 기능안전 관리, PSM & AI 안전성 평가, SIL 적합성 검토 \\
품질 관리부 & 품질 인증, 검사 관리 & AI 품질 검사 기준 관리, 골든 샘플 \\
정보보안팀 & IT 보안, 접근 제어 & OT-IT 융합 보안, IEC 62443 준수 \\
생산 기술팀 & 공정 최적화, 설비 관리 & 예지보전 AI 운영, 드리프트 모니터링 \\
데이터팀/AI팀 & 모델 개발, 데이터 관리 & 모델 인벤토리, 검증, 모니터링 \\
경영진 & 전략 결정, 투자 승인 & AI 거버넌스 정책 승인, 리스크 감독 \\
\bottomrule
\end{tabularx}
\end{table}

\begin{practicaltool}{IEC 61508 AI 적용 가이드 요약}
\textbf{원칙 1}: AI는 SIL 2 이상의 안전 기능에서 단독 제어 역할을 수행해서는 안 된다.\\
\textbf{원칙 2}: AI의 역할은 보조(Advisory) 또는 감시(Monitoring)에 한정하고, 최종 안전 기능은 검증된 하드웨어가 담당한다.\\
\textbf{원칙 3}: AI 모델의 ODD(운영 설계 영역)를 명확히 정의하고, ODD 이탈 시 안전한 대체 행동을 보장한다.\\
\textbf{원칙 4}: AI 모델의 성능 드리프트를 지속적으로 모니터링하며, 임계값 초과 시 자동 대응한다.\\
\textbf{원칙 5}: AI 시스템의 모든 판정과 제어 명령은 추적 가능하게 기록하며, 최소 5년간 보관한다.\\
\textbf{원칙 6}: 정기적 재검증(SIL 등급에 따라 분기~연간)을 수행하고, 결과를 문서화한다.\\
\textbf{원칙 7}: OT-IT 융합 환경에서 보안 거버넌스(IEC 62443)를 안전 거버넌스(IEC 61508)와 통합 관리한다.
\end{practicaltool}

%----------------------------------------------------------
\section{산업 AI 모델 변경 관리(Change Management)}
\label{sec:19.6}

산업 AI에서 모델 변경은 공정 변경과 동일한 수준의 관리가 필요하다. 특히 안전 관련 AI 시스템의 경우, 모델 재학습이나 파라미터 변경이 예상치 못한 동작 변화를 초래할 수 있어 엄격한 변경 관리(Change Management) 프로세스가 요구된다.

\begin{table}[htbp]
\centering
\caption{산업 AI 모델 변경 등급과 관리 절차}
\label{tab:19_change}
\begin{tabularx}{\textwidth}{p{2cm}p{3cm}Xp{3cm}}
\toprule
\textbf{변경 등급} & \textbf{변경 유형} & \textbf{관리 절차} & \textbf{승인 권한} \\
\midrule
Class A & 안전 관련 모델 변경, 알고리즘 교체 & HAZOP 재수행, 전체 검증, 위험성 평가 & 안전관리실장 + CTO \\
Class B & 성능 개선 재학습, 하이퍼파라미터 변경 & 골든 샘플 테스트, 병렬 운영 검증 & 기술팀장 \\
Class C & 임계값 미세 조정, 시각화 변경 & 단위 테스트, 변경 기록 & 담당 엔지니어 \\
Class D & 로그/보고 형식 변경 & 변경 기록만 & 담당자 자체 승인 \\
\bottomrule
\end{tabularx}
\end{table}

모든 변경은 변경 전 상태의 백업, 변경 사유 문서화, 변경 후 검증 결과 기록, 롤백(Rollback) 절차 확인을 포함해야 한다. Class A 변경의 경우 변경 시행 전 최소 2주간의 사전 검토 기간과 관련 부서 협의가 필수적이다.

\paragraph{모델 버전 관리와 롤백}

산업 AI 모델의 버전 관리는 소프트웨어 버전 관리 이상의 엄격함이 요구된다. 모든 프로덕션 모델은 최소 3세대의 이전 버전을 즉시 배포 가능한 상태로 유지해야 하며, 롤백 시 데이터 정합성과 운영 연속성이 보장되어야 한다.

\begin{table}[htbp]
\centering
\caption{산업 AI 모델 버전 관리 요건}
\label{tab:19_version}
\begin{tabularx}{\textwidth}{p{3cm}Xp{3.5cm}}
\toprule
\textbf{관리 항목} & \textbf{내용} & \textbf{기준} \\
\midrule
버전 보관 & 프로덕션 모델 이전 버전 보관 & 최소 3세대, 최대 1년 \\
롤백 시간 & 이전 버전으로 복원 소요 시간 & 5분 이내 (자동 롤백) \\
A/B 테스트 & 신모델 병렬 운영 & 최소 2주, 성능 비교 \\
모델 카드 & 모델별 메타데이터 문서 & 학습 데이터, 성능, 한계 기록 \\
감사 이력 & 모든 배포/롤백 이벤트 기록 & 영구 보관 \\
\bottomrule
\end{tabularx}
\end{table}

%----------------------------------------------------------
\subsection*{제19장 요약}

본 장에서는 제조 및 산업 현장의 AI 거버넌스를 다루었다. 물리적 안전과 직결되는 특성상 IEC 61508(SIL 1-4), ISO 13849(PL a-e), UL 4600 등 기능안전 표준의 준수가 필수적이며, AI의 본질적 한계(비결정론성, 블랙박스 특성)로 인해 SIL 2 이상에서는 AI 단독 제어가 허용되지 않는다는 점을 확인하였다. OT-IT 융합 보안(퍼듀 모델, IEC 62443, 데이터 다이오드)을 통한 사이버 보안 거버넌스와 예지보전 AI의 센서 데이터 관리, 드리프트 탐지, 알림 체계를 상세 코드와 함께 제시하였다. 품질검사 AI의 검사 기준 일관성 관리와 추적 체계, 자율로봇/AMR의 안전 영역 관리와 인간-로봇 협업 거버넌스, 디지털 트윈의 가상-물리 일치성 검증을 다루었으며, H제철과 S전자의 실제 사례를 통해 물리적 가드레일과 드리프트 관리의 실질적 구현 방안을 검토하였다.
